\section*{Candidate Titles}

\noindent\textbf{1. Recommended: Enforcement Activation and Exit Foreclosure in Rule-Governed Organizations}\\
\textit{Justification:} This title states the paper's central separation directly, uses organization-theory language, and leaves religion as the empirical domain.

\noindent\textbf{2. Screening or Foreclosure? The Club-Goods Puzzle of Organizational Strictness}\\
\textit{Justification:} This title foregrounds the identification problem created when the same requirements both select committed members and impede departure.

\noindent\textbf{3. Two Sources of Exit Foreclosure: External Closure, Internal Absorption, and Generational Replacement}\\
\textit{Justification:} This title makes the two structural routes to foreclosure visible while naming the processes that constitute the internal route.

\noindent\textbf{4. Activation, Foreclosure, and Privilege: A Structural Decomposition of Organizational Capture}\\
\textit{Justification:} This title foregrounds the decomposition of capture into enforcement activation, exit conditions, and the institutional production of concentrated punishment.

\noindent\textbf{5. When Enforcement Holds: Code Observability, Sanctioning Rewards, and Foreclosed Exit}\\
\textit{Justification:} This title highlights the joint activation frontier and distinguishes an active apparatus from one that durably retains its subjects.

\section*{Abstract Variants}

\subsection*{Variant A: Puzzle-Led}

Strictness poses a club-goods problem: demanding requirements can screen committed members, yet the same requirements can foreclose exit, so observational data cannot separate selection from lock-in. An agent-based model supplies the instrument by varying code observability, enforcement reward, institutional privilege, and exit capacity independently. With exit open, 0 of 1,890 runs reach capture despite an active rate of 0.598. Under imposed closure, capture rises to 0.889 at closure 0.95, but the low-observability, low-reward cell remains at 0 capture at every closure level. Capture therefore combines an activated enforcement apparatus with foreclosed exit, while activation depends jointly on observability and reward. Privilege also manufactures 84 percent of top-5-percent punishment concentration. The mechanism is content-neutral, and no claim is made that codified belief tends toward coercion.

\noindent\textit{Word count: 124 words. Within the 100--150 word limit: yes.}

\subsection*{Variant B: Findings-Led}

Organizational capture decomposes into enforcement activation, exit foreclosure, and institutionally allocated punishment concentration. Three results follow. First, activation is governed jointly by code observability and enforcement reward: with exit open, the active rate is 0.598, yet 0 of 1,890 runs reach capture. Second, exit closure is necessary but insufficient. Capture rises to 0.889 at imposed closure 0.95, while the low-observability, low-reward cell records 0 capture at every closure level. Third, punishment concentration is manufactured by privilege: the top-5-percent active share rises from a 0.091 floor to a 0.580 ceiling, with 84 percent of ceiling concentration attributable to the privilege bundle. Exit foreclosure can arise through external imposition or internal absorption and generational replacement. The mechanism is content-neutral, and no claim is made that codified belief tends toward coercion.

\noindent\textit{Word count: 128 words. Within the 100--150 word limit: yes.}

\subsection*{Variant C: Organization-Theory-Led}

Delegated sanctioning authority can make rule-governed organizations durable only when an enforcement apparatus activates and members' exit is foreclosed. Religion is the empirical domain for an agent-based model that separates code observability, enforcement reward, institutional privilege, and exit capacity. Enforcement remains active in 0.598 of open-exit runs, but 0 of 1,890 reach capture. Imposed closure raises capture to 0.889 at closure 0.95, except that the low-observability, low-reward cell remains at 0 capture throughout. Sanctioning authority also structures who punishes: the top-5-percent active share is 0.091 without privilege and 0.580 with the full privilege bundle, making 84 percent of ceiling concentration privilege-manufactured. Exit foreclosure arises through external imposition or internal absorption and generational replacement. The mechanism is content-neutral, and no claim is made that codified belief tends toward coercion.

\noindent\textit{Word count: 128 words. Within the 100--150 word limit: yes.}

\section*{Keywords}

organizational capture, exit foreclosure, delegated sanctioning, organizational strictness, institutional privilege, agent-based modeling

\bigskip
\noindent ABSTRACT AND TITLE DONE
